\chapter{Introduction}

\section{Research Background}
Artificial intelligence (AI) is reshaping the interfaces through which brands and consumers meet in digital environments. Across social platforms, e-commerce systems, search interfaces, chatbots, and content feeds, AI increasingly helps determine how messages are produced, personalised, delivered, and interpreted. This matters not only because AI changes marketing efficiency, but because it changes the communicative form of digital brand interaction itself. Recent scholarship suggests that AI is no longer a peripheral back-end tool used only for optimisation; it increasingly participates in the design of consumer-facing communication, interaction, and media experience \cite{haleem2022ai,kumar2024aipowered,jain2024twodecades}.

Recent industry evidence supports the scale of this transition. McKinsey reports that 88\% of organisations now use AI regularly in at least one business function, with marketing and sales among the most common deployment areas, yet most firms remain in experimentation or piloting rather than full-scale transformation \cite{mckinsey2025stateai}. In parallel, Deloitte reports that GenAI is already being used in retail and consumer products for marketing and branding, audience segmentation, copy generation, sentiment analysis, and hyper-personalised shopping support, even though most organisations in these sectors remain in the early stages of adoption \cite{deloitte2025retailgenai}. These figures are useful not as the main frame of the argument, but as evidence that AI-mediated communication has become a practical and visible feature of contemporary digital media environments.

The rapid diffusion of generative artificial intelligence (GenAI) has intensified this transformation. In contrast to earlier forms of predictive AI, which primarily supported classification, targeting, and forecasting, GenAI can directly produce marketing content such as text, images, product descriptions, promotional copy, and interactive brand messages. As a result, AI is no longer only involved in analysing consumer data; it is also participating in the creation of the marketing message itself \cite{hermann2024consumer,ubal2024framework}. This development raises new questions about whether consumers interpret AI-mediated communication in the same way as human-created communication.

From an interactive digital media perspective, this shift matters because consumer responses are not determined solely by functional efficiency. Consumers also evaluate whether an AI-mediated message, interface, or interaction seems trustworthy, authentic, relevant, credible, intrusive, or helpful. These judgments shape broader outcomes such as attitudes towards the brand, willingness to engage, and purchase intention. Although AI can improve convenience and personalisation, its use may simultaneously generate unease when consumers perceive automation as impersonal, deceptive, or overly dependent on personal data \cite{kumar2024aipowered,mohdamin2025adoption}. Understanding consumer perception is therefore central to understanding not only AI adoption in marketing, but also the changing character of digital communication and mediated interaction.

\section{Research Problem}
The use of AI in marketing presents a tension between commercial opportunity and consumer acceptance. On the one hand, AI promises improved efficiency, scalability, and relevance. Personalised recommendations, automated customer support, and AI-generated promotional content can make marketing more responsive to consumer needs and organisational goals. On the other hand, these same practices may trigger concerns about privacy, transparency, manipulation, authenticity, and the erosion of human connection in brand communication \cite{ubal2024framework,hermann2024consumer}.

This tension becomes particularly visible in studies of AI-generated and AI-disclosed marketing content. Emerging empirical work shows that consumers often respond differently when they know a message, image, or advertisement has been created by AI rather than by a human. For example, emotional marketing communication attributed to AI can reduce perceived authenticity and weaken favourable consumer responses \cite{kirk2025aiauthorship}. Similarly, disclosure labels can increase novelty while also decreasing authenticity, producing mixed effects on attitudes and purchase intention \cite{shi2026disclosure}. Studies on AI-generated advertising and visual content also suggest that trust and purchase intention may decline when AI authorship is salient, especially in high-involvement or emotionally sensitive contexts \cite{israfilzade2025advertising,belanche2025images,zhang2025images}.

At the same time, not all findings are negative. Research on AI-powered personalisation indicates that relevance, usefulness, and trust can mediate positive consumer responses, particularly when AI applications are perceived as valuable rather than merely automated \cite{an2025personalized,srivastava2026social}. This suggests that consumer perception of AI marketing is neither uniformly positive nor uniformly negative. Instead, it appears to depend on a set of interconnected perceptual dimensions that operate differently across contexts, content types, and degrees of AI visibility.

Despite growing interest in this area, the literature remains fragmented. Existing studies often examine trust, authenticity, privacy, usefulness, or purchase intention as separate issues, and many focus on specific applications such as AI-generated advertising, recommendation systems, or disclosure labels. As a result, there is still a need for a more integrated account of how AI and GenAI influence consumer perception across marketing contexts. This dissertation addresses that need by synthesising the literature through a thematic review focused explicitly on the consumer perspective.

\section{Research Aim}
The aim of this dissertation is to examine how artificial intelligence and generative artificial intelligence shape consumer perception in AI-mediated communication across digital marketing contexts.

\section{Research Objectives}
\begin{enumerate}
    \item To review academic literature on AI and generative AI in digital marketing and brand communication.
    \item To identify the main consumer perception themes that recur across this literature.
    \item To analyse the mechanisms through which trust, authenticity, personalisation, privacy concerns, and perceived usefulness shape consumer responses to AI-mediated communication.
    \item To explain the conditions under which AI-mediated communication is more likely to generate favourable or unfavourable consumer responses.
\end{enumerate}

\section{Research Questions}
\begin{enumerate}
    \item What are the main dimensions through which AI and generative AI shape consumer perception in digital marketing communication?
    \item Through what mechanisms do trust, authenticity, personalisation, privacy concerns, and perceived usefulness influence consumer responses to AI-mediated communication?
    \item Under what conditions are consumer responses to AI and generative AI in digital marketing communication more likely to become positive or negative?
\end{enumerate}

\section{Significance of the Study}
This study is significant for both academic and practical reasons. Academically, it contributes to a developing body of literature at the intersection of AI, consumer behaviour, digital communication, and interactive media. While research on AI in marketing has expanded rapidly, the consumer perspective remains conceptually dispersed across different applications and theoretical lenses \cite{jain2024twodecades,ubal2024framework}. By bringing together literature on predictive AI, generative AI, advertising, personalisation, disclosure, and mediated interaction, this dissertation offers a clearer account of the perceptual mechanisms that shape consumer responses.

Practically, the study is relevant for marketers and organisations that are increasingly incorporating AI into customer-facing activities. Businesses may assume that efficiency and innovation automatically improve consumer evaluations, yet the literature suggests that acceptance depends on how consumers interpret the meaning of AI use. If AI-generated content is seen as useful and relevant, it may strengthen engagement and purchase intention; if it is seen as inauthentic, misleading, or invasive, it may damage trust and brand relationships \cite{kirk2025aiauthorship,an2025personalized,israfilzade2025advertising}. A better understanding of these dynamics can help organisations deploy AI more responsibly and strategically.

The study is also relevant to interactive digital media design and digital communication practice. As AI becomes embedded in chat interfaces, recommendation systems, disclosure cues, automated service encounters, and generative content workflows, design decisions increasingly shape how AI is experienced by users. The dissertation therefore has implications for how AI-mediated brand interactions are framed, how disclosure is designed at the user experience level, and how organisations position GenAI content within broader digital media production and communication processes.

\section{Scope of the Study}
This dissertation focuses on consumer perception in digital marketing communication rather than on the technical performance of AI systems. It examines literature on AI-enabled practices including personalised advertising, recommendation systems, AI-generated advertising content, chatbots, and social media communication. The study considers both traditional predictive AI and newer forms of generative AI, especially where AI becomes visible in the marketing message or consumer experience.

More specifically, the dissertation is concerned with consumer-AI interaction in digital media environments rather than with organisational adoption at the level of marketing strategy alone.

The dissertation is limited to secondary research and does not involve primary data collection. Its purpose is not to test causal relationships empirically, but to synthesise and interpret existing academic evidence. The emphasis is therefore on identifying recurring themes, tensions, and patterns in the literature rather than on evaluating the engineering performance of AI technologies.

\section{Dissertation Structure}
This dissertation is organised into five chapters. Chapter 1 introduces the research background, problem, aim, objectives, research questions, and scope of the study. Chapter 2 reviews the literature on AI, generative AI, and consumer perception in digital marketing communication, and develops the conceptual basis for the study. Chapter 3 explains the methodology, including the literature selection process and thematic analysis approach. Chapter 4 presents and discusses the main themes identified from the reviewed literature. Chapter 5 concludes the dissertation by summarising the findings, outlining the implications of the study, and suggesting directions for future research.
